%TEX program=xelatex
\documentclass[zihao=false,linespread=1.5,heading=true,a4paper,oneside,fontset=windows]{ctexbook}
\usepackage[fontsize=12pt]{fontsize}
\usepackage[utf8]{inputenc}
\pagestyle{plain}
% packages for page layout
\usepackage{geometry}
\geometry{
a4paper,
total={171.8mm,246.2mm},
left=19.1mm,
top=25.4mm,
}

% package for contents
\usepackage{tocbibind} 


\usepackage{afterpage}
% package for colors
\usepackage{xcolor}
\usepackage{color}
\definecolor{coverbackground}{RGB}{128,128,128}%灰色
% Colors
\definecolor{BLUE_E}{HTML}{1C758A}
\definecolor{BLUE_D}{HTML}{29ABCA}
\definecolor{BLUE_C}{HTML}{58C4DD}
\definecolor{BLUE_B}{HTML}{9CDCEB}
\definecolor{BLUE_A}{HTML}{C7E9F1}
\definecolor{TEAL_E}{HTML}{49A88F}
\definecolor{TEAL_D}{HTML}{55C1A7}
\definecolor{TEAL_C}{HTML}{5CD0B3}
\definecolor{TEAL_B}{HTML}{76DDC0}
\definecolor{TEAL_A}{HTML}{ACEAD7}
\definecolor{GREEN_E}{HTML}{699C52}
\definecolor{GREEN_D}{HTML}{77B05D}
\definecolor{GREEN_C}{HTML}{83C167}
\definecolor{GREEN_B}{HTML}{A6CF8C}
\definecolor{GREEN_A}{HTML}{C9E2AE}
\definecolor{YELLOW_E}{HTML}{E8C11C}
\definecolor{YELLOW_D}{HTML}{F4D345}
\definecolor{YELLOW_C}{HTML}{FFFF00}
\definecolor{YELLOW_B}{HTML}{FFEA94}
\definecolor{YELLOW_A}{HTML}{FFF1B6}
\definecolor{GOLD_E}{HTML}{C78D46}
\definecolor{GOLD_D}{HTML}{E1A158}
\definecolor{GOLD_C}{HTML}{F0AC5F}
\definecolor{GOLD_B}{HTML}{F9B775}
\definecolor{GOLD_A}{HTML}{F7C797}
\definecolor{RED_E}{HTML}{CF5044}
\definecolor{RED_D}{HTML}{E65A4C}
\definecolor{RED_C}{HTML}{FC6255}
\definecolor{RED_B}{HTML}{FF8080}
\definecolor{RED_A}{HTML}{F7A1A3}
\definecolor{MAROON_E}{HTML}{94424F}
\definecolor{MAROON_D}{HTML}{A24D61}
\definecolor{MAROON_C}{HTML}{C55F73}
\definecolor{MAROON_B}{HTML}{EC92AB}
\definecolor{MAROON_A}{HTML}{ECABC1}
\definecolor{PURPLE_E}{HTML}{644172}
\definecolor{PURPLE_D}{HTML}{715582}
\definecolor{PURPLE_C}{HTML}{9A72AC}
\definecolor{PURPLE_B}{HTML}{B189C6}
\definecolor{PURPLE_A}{HTML}{CAA3E8}
\definecolor{GREY_E}{HTML}{222222}
\definecolor{GREY_D}{HTML}{444444}
\definecolor{GREY_C}{HTML}{888888}
\definecolor{GREY_B}{HTML}{BBBBBB}
\definecolor{GREY_A}{HTML}{DDDDDD}
\definecolor{WHITE}{HTML}{FFFFFF}
\definecolor{BLACK}{HTML}{000000}
\definecolor{GREY_BROWN}{HTML}{736357}
\definecolor{DARK_BROWN}{HTML}{8B4513}
\definecolor{LIGHT_BROWN}{HTML}{CD853F}
\definecolor{PINK}{HTML}{D147BD}
\definecolor{LIGHT_PINK}{HTML}{DC75CD}
\definecolor{GREEN_SCREEN}{HTML}{00FF00}
\definecolor{ORANGE}{HTML}{FF862F}

%COLORMAP_3B1B [BLUE_E, GREEN, YELLOW, RED]

% packages for fonts
\usepackage{amsmath}

\usepackage{fontspec}
\usepackage{amsfonts}
\usepackage{mathrsfs}
\usepackage{mathtools}
\usepackage{amssymb}
\usepackage{bm}
\DeclareFontFamily{U}{rsfs}{\skewchar\font127}
\DeclareFontShape{U}{rsfs}{m}{n}{%
   <-6> rsfs5
   <6-8> rsfs7
   <8-> rsfs10
}{}   

% package for tikz
\usepackage{tcolorbox}
\usepackage{setspace}
\tcbuselibrary{skins, breakable, theorems}
\tcbuselibrary{most}
\usepackage{framed}
\usepackage{indentfirst}
\usepackage{tikz}
\usetikzlibrary{arrows.meta}
\usetikzlibrary{fpu}
\usetikzlibrary{math}
\usetikzlibrary{cd}

% packages for graphics
\usepackage{zhlipsum}
\usepackage{graphicx}

\graphicspath{ {./images/}{./images/} }
\usepackage{wrapfig}
\usepackage{subcaption}
\usepackage{capt-of}
\usepackage{cutwin}

\usepackage{fancyhdr}
\pagestyle{fancy}
\fancyhead{}
\fancyhead[L]{\leftmark}
\fancyfoot{}
\fancyfoot[C]{\thepage}

% package for quotes
\usepackage{csquotes}

% redefine itemize
\usepackage{enumitem}% http://ctan.org/pkg/enumitem

% pacakage of contents
\usepackage{lipsum}
\usepackage{titletoc}

% packages for footnotes
\renewcommand{\thefootnote}{\fnsymbol{footnote}}
\usepackage{perpage}
\MakePerPage{footnote} 

% package for svg
\usepackage{svg} % 需使用包

% packages for definitions, theorems, proofs.
\usepackage{amsthm}

%\oiint problem
\usepackage{esint}
% counter controlling
\usepackage{chngcntr}

\usepackage{datetime2}
% package for hyperlinks
\usepackage{multicol}
\usepackage{tabularx}
\usepackage{makeidx}
\makeindex
\renewcommand{\indexname}{公式与主题索引}
\usepackage{hyperref}
\hypersetup{citecolor=magenta,linkcolor=black,urlcolor=blue,hidelinks} % 设置引用链接颜色

%定制章节编号格式和计数器
\renewcommand\thesection{\arabic{section}}

% package for subfiles
\usepackage{subfiles}
%=====================================================================================================================
\ctexset{part={name={第,篇},number=\chinese{part}}}
\begin{document}

\thispagestyle{empty}%目录页不显示页码
\hypersetup{pageanchor=false}
\subfile{cover/cover.tex}

\setcounter{page}{1}%从下面开始编页码
\hypersetup{pageanchor=true}
\subfile{preface/preface.tex}
\newpage

\subfile{introduction/introduction.tex}
\newpage

\tableofcontents
\newpage

% 按主题分篇，并在各分支内由基础到进阶排列。
\part{基础数学}\label{part:basic-mathematics}

\chapter{初等数学}\label{chap:elementary-mathematics}
\subfile{chapter/初等数学.tex}

\chapter{平面几何}\label{chap:plane-geometry}
\subfile{chapter/平面几何.tex}

\chapter{立体几何}\label{chap:solid-geometry}
\subfile{chapter/立体几何.tex}

\chapter{初等数论}\label{chap:elementary-number-theory}
\subfile{chapter/初等数论.tex}

\chapter{高等代数}\label{chap:linear-algebra}
\subfile{chapter/高等代数.tex}

\chapter{数学分析}\label{chap:mathematical-analysis}
\subfile{chapter/数学分析.tex}

\part{逻辑、集合与离散数学}\label{part:logic-discrete}

\chapter{数理逻辑}\label{chap:mathematical-logic}
\subfile{chapter/数理逻辑.tex}

\chapter{集合论}\label{chap:set-theory}
\subfile{chapter/集合论.tex}

\chapter{离散数学}\label{chap:discrete-mathematics}
\subfile{chapter/离散数学.tex}

\chapter{组合数学与图论}\label{chap:combinatorics-and-graph-theory}
\subfile{chapter/组合数学与图论.tex}

\part{代数结构}\label{part:algebraic-structures}

\chapter{群论}\label{chap:group-theory}
\subfile{chapter/群论.tex}

\chapter{环论}\label{chap:ring-theory}
\subfile{chapter/环论.tex}

\chapter{域论}\label{chap:field-theory}
\subfile{chapter/域论.tex}

\chapter{模论}\label{chap:module-theory}
\subfile{chapter/模论.tex}

\part{拓扑、测度与分析}\label{part:analysis-foundations}

\chapter{一般拓扑}\label{chap:general-topology}
\subfile{chapter/一般拓扑.tex}

\chapter{测度论}\label{chap:measure-theory}
\subfile{chapter/测度论.tex}

\chapter{实分析}\label{chap:real-analysis}
\subfile{chapter/实分析.tex}

\chapter{复分析}\label{chap:complex-analysis}
\subfile{chapter/复分析.tex}

\chapter{泛函分析}\label{chap:functional-analysis}
\subfile{chapter/泛函分析.tex}

\chapter{凸分析}\label{chap:convex-analysis}
\subfile{chapter/凸分析.tex}

\chapter{调和分析}\label{chap:harmonic-analysis}
\subfile{chapter/调和分析.tex}

\part{方程、变分与动力系统}\label{part:equations-variations}

\chapter{常微分方程}\label{chap:ordinary-differential-equations}
\subfile{chapter/常微分方程.tex}

\chapter{数值分析}\label{chap:numerical-analysis}
\subfile{chapter/数值分析.tex}

\chapter{变分法}\label{chap:calculus-of-variations}
\subfile{chapter/变分法.tex}

\chapter{偏微分方程}\label{chap:partial-differential-equations}
\subfile{chapter/偏微分方程.tex}

\chapter{动力系统}\label{chap:dynamical-systems}
\subfile{chapter/动力系统.tex}

\chapter{分形分析}\label{chap:fractal-analysis}
\subfile{chapter/分形分析.tex}

\part{概率、统计与随机过程}\label{part:probability-statistics}

\chapter{概率论}\label{chap:probability}
\subfile{chapter/概率论.tex}

\chapter{数理统计}\label{chap:mathematical-statistics}
\subfile{chapter/数理统计.tex}

\chapter{随机过程}\label{chap:stochastic-processes}
\subfile{chapter/随机过程.tex}

\part{高阶几何与拓扑}\label{part:advanced-geometry}

\chapter{微分几何}\label{chap:differential-geometry}
\subfile{chapter/微分几何.tex}

\chapter{代数拓扑}\label{chap:algebraic-topology}
\subfile{chapter/代数拓扑.tex}

\chapter{微分拓扑}\label{chap:differential-topology}
\subfile{chapter/微分拓扑.tex}

\chapter{代数几何}\label{chap:algebraic-geometry}
\subfile{chapter/代数几何.tex}

\part{进阶数论}\label{part:advanced-number-theory}

\chapter{几何数论}\label{chap:geometry-of-numbers}
\subfile{chapter/几何数论.tex}

\chapter{解析数论}\label{chap:analytic-number-theory}
\subfile{chapter/解析数论.tex}

\chapter{代数数论}\label{chap:algebraic-number-theory}
\subfile{chapter/代数数论.tex}


\clearpage
\begin{thebibliography}{99}
\label{back:references}
\item[] 以下列出各章的基础教材、作者讲义与本轮补充的主要核对来源。章节开头的引用提供阅读入口，新增专题在节内注明相应来源；编号与正文引用对应。在线资源访问日期：2026年11月11日。

\bibitem{openstax} Jay Abramson et al.. \emph{Algebra and Trigonometry}. 2nd ed., OpenStax, 2021. \href{https://openstax.org/details/books/algebra-and-trigonometry-2e}{出版或作者资源链接}.

\bibitem{euclid} Euclid; Richard Fitzpatrick (ed. and trans.). \emph{Euclid's Elements of Geometry}. revised edition, 2008. \href{https://farside.ph.utexas.edu/Books/Euclid/Elements.pdf}{出版或作者资源链接}.

\bibitem{geometry} Anna Felikson. \emph{Geometry: Lecture Notes}. Durham University, 2021. \href{https://www.maths.dur.ac.uk/users/anna.felikson/Geometry/Geometry21/notes.pdf}{出版或作者资源链接}.

\bibitem{hlp} G. H. Hardy, J. E. Littlewood and G. Pólya. \emph{Inequalities}. 2nd ed., Cambridge University Press, 1952. \href{https://assets.cambridge.org/97805213/58804/toc/9780521358804_toc.pdf}{出版或作者资源链接}.

\bibitem{shoup} Victor Shoup. \emph{A Computational Introduction to Number Theory and Algebra}. 2nd ed., author's electronic edition. \href{https://www.shoup.net/ntb/}{出版或作者资源链接}.

\bibitem{lte} Maya Kądziołka. \emph{Lifting the Exponent}. Archive of Formal Proofs, 2021. \href{https://www.isa-afp.org/entries/Lifting_the_Exponent.html}{出版或作者资源链接}.

\bibitem{matrix} Kaare Brandt Petersen and Michael Syskind Pedersen. \emph{The Matrix Cookbook}. 2012-11-15. \href{https://www.math.uwaterloo.ca/~hwolkowi/matrixcookbook.pdf}{出版或作者资源链接}.

\bibitem{dlmf} NIST. \emph{Digital Library of Mathematical Functions}. online reference, version 1.2.8, 2026. \href{https://dlmf.nist.gov/}{出版或作者资源链接}.

\bibitem{freund} Anton Freund. \emph{An Introduction to Mathematical Logic}. arXiv:2310.09921v2, 2023. \href{https://arxiv.org/abs/2310.09921}{出版或作者资源链接}.

\bibitem{weiss} William A. R. Weiss. \emph{An Introduction to Set Theory}. University of Toronto, 2008. \href{https://www.math.toronto.edu/weiss/set_theory.pdf}{出版或作者资源链接}.

\bibitem{jech} Thomas Jech. \emph{Set Theory}. Third Millennium Edition, Springer, 2003. \href{https://doi.org/10.1007/3-540-44761-X}{出版或作者资源链接}.

\bibitem{mcs} Eric Lehman, F. Thomson Leighton and Albert R. Meyer. \emph{Mathematics for Computer Science}. MIT OpenCourseWare, 2015. \href{https://www.ocw.mit.edu/courses/6-042j-mathematics-for-computer-science-spring-2015/resources/mit6_042js15_textbook_pdf/}{出版或作者资源链接}.

\bibitem{cameron} Peter J. Cameron. \emph{Combinatorics: Topics, Techniques, Algorithms}. Cambridge University Press, 1994. \href{https://webspace.maths.qmul.ac.uk/p.j.cameron/comb/}{出版或作者资源链接}.

\bibitem{milneGT} J. S. Milne. \emph{Group Theory}. version 4.01, 2025. \href{https://www.jmilne.org/math/CourseNotes/GT.pdf}{出版或作者资源链接}.

\bibitem{milneFT} J. S. Milne. \emph{Fields and Galois Theory}. version 5.10, 2022. \href{https://www.jmilne.org/math/CourseNotes/FT.pdf}{出版或作者资源链接}.

\bibitem{milneCA} J. S. Milne. \emph{A Primer of Commutative Algebra}. version 4.03, 2020. \href{https://www.jmilne.org/math/xnotes/ca.html}{出版或作者资源链接}.

\bibitem{stacks} The Stacks Project Authors. \emph{The Stacks Project}. online reference; algebra, modules, schemes and cohomology. \href{https://stacks.math.columbia.edu/}{出版或作者资源链接}.

\bibitem{topology} John M. Lee. \emph{Introduction to Topological Manifolds}. 2nd ed., Springer, 2011. \href{https://sites.math.washington.edu/~lee/Books/ITM/}{出版或作者资源链接}.

\bibitem{brezis} Haim Brezis. \emph{Functional Analysis, Sobolev Spaces and Partial Differential Equations}. Springer, 2011. \href{https://doi.org/10.1007/978-0-387-70914-7}{出版或作者资源链接}.

\bibitem{bhatia} Rajendra Bhatia. \emph{Matrix Analysis}. Springer, 1997. \href{https://doi.org/10.1007/978-1-4612-0653-8}{出版或作者资源链接}.

\bibitem{kinnunen} Juha Kinnunen. \emph{Sobolev Spaces}. Aalto University, lecture notes, 2026. \href{https://math.aalto.fi/~jkkinnun/files/sobolev_spaces.pdf}{出版或作者资源链接}.

\bibitem{villani} Cédric Villani. \emph{Optimal Transport: Old and New}. Springer, 2009. \href{https://doi.org/10.1007/978-3-540-71050-9}{出版或作者资源链接}.

\bibitem{hunter} John K. Hunter. \emph{Notes on Partial Differential Equations}. University of California, Davis, 2014. \href{https://math.ucdavis.edu/~hunter/pdes/pde_notes.pdf}{出版或作者资源链接}.

\bibitem{complex} Elias M. Stein and Rami Shakarchi. \emph{Complex Analysis}. Princeton University Press, 2003. \href{https://press.princeton.edu/books/hardcover/9780691113852/complex-analysis}{出版或作者资源链接}.

\bibitem{convex} R. Tyrrell Rockafellar. \emph{Convex Analysis}. Princeton University Press, 1970. \href{https://press.princeton.edu/books/paperback/9780691015866/convex-analysis}{出版或作者资源链接}.

\bibitem{chambolle} Antonin Chambolle. \emph{Continuous Optimization, an Introduction}. lecture notes, 2024. \href{https://www.ceremade.dauphine.fr/~chambolle/data/medias/optimisation/coursopti2024.pdf}{出版或作者资源链接}.

\bibitem{boyd} Stephen Boyd and Lieven Vandenberghe. \emph{Convex Optimization}. Cambridge University Press, 2004. \href{https://web.stanford.edu/~boyd/cvxbook/}{出版或作者资源链接}.

\bibitem{sion} Maurice Sion. \emph{On General Minimax Theorems}. Pacific Journal of Mathematics 8 (1958), 171--176. \href{https://doi.org/10.2140/pjm.1958.8.171}{出版或作者资源链接}.

\bibitem{harmonic} Terence Tao. \emph{Lecture Notes for 247B: Paraproducts}. UCLA, 2007, Notes 6. \href{https://www.math.ucla.edu/~tao/247b.1.07w/notes6.pdf}{出版或作者资源链接}.

\bibitem{williams} Mark Williams. \emph{Notes on Harmonic Analysis}. lecture notes, 2022. \href{https://markwilliams.web.unc.edu/wp-content/uploads/sites/19674/2022/01/notesonharmonicanalysisB.pdf}{出版或作者资源链接}.

\bibitem{beckner} William Beckner. \emph{Inequalities in Fourier Analysis}. Annals of Mathematics 102 (1975), 159--182. \href{https://annals.math.princeton.edu/1975/102-1/p11}{出版或作者资源链接}.

\bibitem{teschl} Gerald Teschl. \emph{Ordinary Differential Equations and Dynamical Systems}. AMS, Graduate Studies in Mathematics 140, 2012. \href{https://www.mat.univie.ac.at/~gerald/ftp/book-ode/}{出版或作者资源链接}.

\bibitem{fnc} Tobin A. Driscoll and Richard J. Braun. \emph{Fundamentals of Numerical Computation}. author's online edition. \href{https://fncbook.com/}{出版或作者资源链接}.

\bibitem{spectral} Lloyd N. Trefethen. \emph{Spectral Methods in MATLAB}. SIAM, 2000. \href{https://doi.org/10.1137/1.9780898719598}{出版或作者资源链接}.

\bibitem{geometric} Ernst Hairer, Christian Lubich and Gerhard Wanner. \emph{Geometric Numerical Integration}. 2nd ed., Springer, 2006. \href{https://doi.org/10.1007/3-540-30666-8}{出版或作者资源链接}.

\bibitem{fem} Alexandre Ern and Jean-Luc Guermond. \emph{Theory and Practice of Finite Elements}. Springer, 2004. \href{https://doi.org/10.1007/978-1-4757-4355-5}{出版或作者资源链接}.

\bibitem{variation} Filip Rindler. \emph{Introduction to the Modern Calculus of Variations}. MA4G6 lecture notes, University of Warwick. \href{https://warwick.ac.uk/fac/sci/maths/people/staff/filip_rindler/cov_ln.pdf}{出版或作者资源链接}.

\bibitem{kinetic} Paul J. Dellar. \emph{Kinetic Theory of Neutral Particles}. MMathPhys course notes, University of Oxford. \href{https://people.maths.ox.ac.uk/dellar/MMPkinetic.html}{出版或作者资源链接}.

\bibitem{falconer} Kenneth Falconer. \emph{Fractal Geometry: Mathematical Foundations and Applications}. 3rd ed., Wiley, 2014. \href{https://www.wiley.com/en-us/fractal-geometry-mathematical-foundations-and-applications-3rd-edition-p-9781118762868}{出版或作者资源链接}.

\bibitem{mtp} Victor Beresnevich and Sanju Velani. \emph{A Mass Transference Principle and the Duffin-Schaeffer Conjecture for Hausdorff Measures}. Annals of Mathematics 164 (2006), 971--992. \href{https://doi.org/10.4007/annals.2006.164.971}{出版或作者资源链接}.

\bibitem{durrett} Rick Durrett. \emph{Probability: Theory and Examples}. 5th ed., author's online edition. \href{https://sites.math.duke.edu/~rtd/PTE/pte.html}{出版或作者资源链接}.

\bibitem{wasserman} Larry Wasserman. \emph{All of Statistics: A Concise Course in Statistical Inference}. Springer, 2004. \href{https://www.stat.cmu.edu/~larry/all-of-statistics/}{出版或作者资源链接}.

\bibitem{legall} Jean-François Le Gall. \emph{Brownian Motion, Martingales, and Stochastic Calculus}. Springer, 2016. \href{https://doi.org/10.1007/978-3-319-31089-3}{出版或作者资源链接}.

\bibitem{lee} John M. Lee. \emph{Introduction to Smooth Manifolds}. 2nd ed., Springer, 2012. \href{https://sites.math.washington.edu/~lee/Books/ISM/}{出版或作者资源链接}.

\bibitem{riemann} John M. Lee. \emph{Introduction to Riemannian Manifolds}. 2nd ed., Springer, 2018. \href{https://sites.math.washington.edu/~lee/Books/RM/}{出版或作者资源链接}.

\bibitem{hatcher} Allen Hatcher. \emph{Algebraic Topology}. Cambridge University Press, 2002. \href{https://pi.math.cornell.edu/~hatcher/AT/ATpage.html}{出版或作者资源链接}.

\bibitem{milneAG} J. S. Milne. \emph{Algebraic Geometry}. version 6.10, 2024. \href{https://www.jmilne.org/math/CourseNotes/ag.html}{出版或作者资源链接}.

\bibitem{cassels} J. W. S. Cassels. \emph{An Introduction to the Geometry of Numbers}. Springer, corrected reprint, 1997. \href{https://doi.org/10.1007/978-3-642-62035-5}{出版或作者资源链接}.

\bibitem{davenport} Harold Davenport; H. L. Montgomery (rev.). \emph{Multiplicative Number Theory}. 3rd ed., Springer, 2000. \href{https://link.springer.com/book/9780387950976}{出版或作者资源链接}.

\bibitem{milneANT} J. S. Milne. \emph{Algebraic Number Theory}. version 3.08, 2020. \href{https://www.jmilne.org/math/CourseNotes/ANT.pdf}{出版或作者资源链接}.

\bibitem{massart} Pascal Massart. \emph{Concentration Inequalities and Model Selection}. Springer, Lecture Notes in Mathematics 1896, 2007. \href{https://doi.org/10.1007/978-3-540-48503-2}{出版或作者资源链接}.

\bibitem{sparsity} Trevor Hastie, Robert Tibshirani and Martin J. Wainwright. \emph{Statistical Learning with Sparsity: The Lasso and Generalizations}. Chapman \& Hall/CRC, 2015. \href{https://wainwrigwork.github.io/publications.html}{出版或作者资源链接}.

\bibitem{topping} Peter M. Topping. \emph{Lectures on the Ricci Flow}. Cambridge University Press, LMS Lecture Note Series 325, 2006. \href{https://warwick.ac.uk/fac/sci/maths/people/staff/peter_topping/rfnotes/}{出版或作者资源链接}.

\bibitem{heatkernels} Nicole Berline, Ezra Getzler and Michèle Vergne. \emph{Heat Kernels and Dirac Operators}. Springer, 1992; Classics in Mathematics reprint, 2004. \href{https://doi.org/10.1007/978-3-642-58088-8}{出版或作者资源链接}.

\bibitem{hatcherK} Allen Hatcher. \emph{Vector Bundles \& K-Theory}. author's online book. \href{https://pi.math.cornell.edu/~hatcher/VBKT/VBpage.html}{出版或作者资源链接}.

\bibitem{symplectic} Ana Cannas da Silva. \emph{Lectures on Symplectic Geometry}. Springer, Lecture Notes in Mathematics 1764, 2001; revised author's text, 2006. \href{https://www.math.ist.utl.pt/~acannas/Books/symplectic.pdf}{出版或作者资源链接}.

\bibitem{toric} David A. Cox, John B. Little and Henry K. Schenck. \emph{Toric Varieties}. AMS, Graduate Studies in Mathematics 124, 2011. \href{https://doi.org/10.1090/gsm/124}{出版或作者资源链接}.

\bibitem{lattices} Daniele Micciancio and Shafi Goldwasser. \emph{Complexity of Lattice Problems: A Cryptographic Perspective}. Springer, 2002. \href{https://doi.org/10.1007/978-1-4615-0897-7}{出版或作者资源链接}.

\bibitem{schneider} Rolf Schneider. \emph{Convex Bodies: The Brunn--Minkowski Theory}. 2nd expanded ed., Cambridge University Press, 2014. \href{https://assets.cambridge.org/97811076/01017/frontmatter/9781107601017_frontmatter.pdf}{出版或作者资源链接}.

\bibitem{tenenbaum} Gérald Tenenbaum; Patrick D. F. Ion (trans.). \emph{Introduction to Analytic and Probabilistic Number Theory}. 3rd ed., AMS, Graduate Studies in Mathematics 163, 2015. \href{https://doi.org/10.1090/gsm/163}{出版或作者资源链接}.

\bibitem{titchmarsh} E. C. Titchmarsh; D. R. Heath-Brown (rev.). \emph{The Theory of the Riemann Zeta-function}. 2nd ed., Oxford University Press, 1986. \href{https://sites.math.rutgers.edu/~zeilberg/EM18/TitchmarshZeta.pdf}{出版或作者资源链接}.

\bibitem{milneCFT} J. S. Milne. \emph{Class Field Theory}. version 4.03, 2020. \href{https://www.jmilne.org/math/CourseNotes/CFT.pdf}{出版或作者资源链接}.

\end{thebibliography}

\clearpage
\phantomsection\label{back:index}
\printindex

\newpage
\thispagestyle{empty}%尾页不显示页码
\subfile{colophon/colophon.tex}
\newpage

\end{document}
